
\documentclass[a4paper,11pt]{article}
\usepackage[a4paper, margin=8em]{geometry}

% usa i pacchetti per la scrittura in italiano
\usepackage[french,italian]{babel}
\usepackage[T1]{fontenc}
\usepackage[utf8]{inputenc}
\frenchspacing 

% usa i pacchetti per la formattazione matematica
\usepackage{amsmath, amssymb, amsthm, amsfonts}

% usa altri pacchetti
\usepackage{gensymb}
\usepackage{hyperref}
\usepackage{standalone}

% imposta il titolo
\title{Appunti Comunicazioni Numeriche}
\author{Luca Seggiani}
\date{2026}

% imposta lo stile
% usa helvetica
\usepackage[scaled]{helvet}
% usa palatino
\usepackage{palatino}
% usa un font monospazio guardabile
\usepackage{lmodern}

\renewcommand{\rmdefault}{ppl}
\renewcommand{\sfdefault}{phv}
\renewcommand{\ttdefault}{lmtt}

% disponi il titolo
\makeatletter
\renewcommand{\maketitle} {
	\begin{center} 
		\begin{minipage}[t]{.8\textwidth}
			\textsf{\huge\bfseries \@title} 
		\end{minipage}%
		\begin{minipage}[t]{.2\textwidth}
			\raggedleft \vspace{-1.65em}
			\textsf{\small \@author} \vfill
			\textsf{\small \@date}
		\end{minipage}
		\par
	\end{center}

	\thispagestyle{empty}
	\pagestyle{fancy}
}
\makeatother

% disponi teoremi
\usepackage{tcolorbox}
\newtcolorbox[auto counter, number within=section]{theorem}[2][]{%
	colback=blue!10, 
	colframe=blue!40!black, 
	sharp corners=northwest,
	fonttitle=\sffamily\bfseries, 
	title=~\thetcbcounter: #2, 
	#1
}

% disponi definizioni
\newtcolorbox[auto counter, number within=section]{definition}[2][]{%
	colback=red!10,
	colframe=red!40!black,
	sharp corners=northwest,
	fonttitle=\sffamily\bfseries,
	title=~\thetcbcounter: #2,
	#1
}

% U.D.M
\newcommand{\amp}{\ensuremath{\mathrm{A}}}
\newcommand{\volt}{\ensuremath{\mathrm{V}}}
\newcommand{\meter}{\ensuremath{\mathrm{m}}}
\newcommand{\second}{\ensuremath{\mathrm{s}}}
\newcommand{\farad}{\ensuremath{\mathrm{F}}}
\newcommand{\henry}{\ensuremath{\mathrm{H}}}
\newcommand{\siemens}{\ensuremath{\mathrm{S}}}

% circuiti
\usepackage{circuitikz}
\usetikzlibrary{babel}

% disegni
\usepackage{pgfplots}
\pgfplotsset{width=10cm,compat=1.9}

% disponi codice
\usepackage{listings}
\usepackage[table]{xcolor}

\lstdefinestyle{codestyle}{
		backgroundcolor=\color{black!5}, 
		commentstyle=\color{codegreen},
		keywordstyle=\bfseries\color{magenta},
		numberstyle=\sffamily\tiny\color{black!60},
		stringstyle=\color{green!50!black},
		basicstyle=\ttfamily\footnotesize,
		breakatwhitespace=false,         
		breaklines=true,                 
		captionpos=b,                    
		keepspaces=true,                 
		numbers=left,                    
		numbersep=5pt,                  
		showspaces=false,                
		showstringspaces=false,
		showtabs=false,                  
		tabsize=2
}

\lstdefinestyle{shellstyle}{
		backgroundcolor=\color{black!5}, 
		basicstyle=\ttfamily\footnotesize\color{black}, 
		commentstyle=\color{black}, 
		keywordstyle=\color{black},
		numberstyle=\color{black!5},
		stringstyle=\color{black}, 
		showspaces=false,
		showstringspaces=false, 
		showtabs=false, 
		tabsize=2, 
		numbers=none, 
		breaklines=true
}

\lstdefinelanguage{javascript}{
	keywords={typeof, new, true, false, catch, function, return, null, catch, switch, var, if, in, while, do, else, case, break},
	keywordstyle=\color{blue}\bfseries,
	ndkeywords={class, export, boolean, throw, implements, import, this},
	ndkeywordstyle=\color{darkgray}\bfseries,
	identifierstyle=\color{black},
	sensitive=false,
	comment=[l]{//},
	morecomment=[s]{/*}{*/},
	commentstyle=\color{purple}\ttfamily,
	stringstyle=\color{red}\ttfamily,
	morestring=[b]',
	morestring=[b]"
}

% disponi sezioni
\usepackage{titlesec}

\titleformat{\section}
	{\sffamily\Large\bfseries} 
	{\thesection}{1em}{} 
\titleformat{\subsection}
	{\sffamily\large\bfseries}   
	{\thesubsection}{1em}{} 
\titleformat{\subsubsection}
	{\sffamily\normalsize\bfseries} 
	{\thesubsubsection}{1em}{}

% disponi alberi
\usepackage{forest}

\forestset{
	rectstyle/.style={
		for tree={rectangle,draw,font=\large\sffamily}
	},
	roundstyle/.style={
		for tree={circle,draw,font=\large}
	}
}

% disponi algoritmi
\usepackage{algorithm}
\usepackage{algorithmic}
\makeatletter
\renewcommand{\ALG@name}{Algoritmo}
\makeatother

% disponi numeri di pagina
\usepackage{fancyhdr}
\fancyhf{} 
\fancyfoot[L]{\sffamily{\thepage}}

\makeatletter
\fancyhead[L]{\raisebox{1ex}[0pt][0pt]{\sffamily{\@title \ \@date}}} 
\fancyhead[R]{\raisebox{1ex}[0pt][0pt]{\sffamily{\@author}}}
\makeatother

\begin{document}
% sezione (data)
\section{Lezione del 10-04-26}

% stili pagina
\thispagestyle{empty}
\pagestyle{fancy}

% testo
\subsection{Sistemi monodimensionali a tempo continuo}
La trasformata di Fourier, vista come modo per trattare i segnali nelle ultime lezioni, può essere utile anche a trattare i \textit{sistemi} fisici (ad esempio, sistemi elettronici, ecc...). 

Un \textbf{sistema} in generale è una trasformazione (un \textit{funzionale} o \textit{operatore}, ecc...) che prende in ingresso un segnale $x(t)$ e lo porta in un secondo segnale $y(t)$, detto uscita. Matematicamente, lo indichiamo come:
$$
y(t) = \mathcal{T} \left[ x(\alpha); t \right]
$$
Ricordiamo che, per una trasformazione $\mathcal{T}$, ad ogni segnale di ingresso $x(t)$ corrisponde un unico segnale di uscita $y(t)$. Può chiaramente essere che il segnale $y(t)$ al tempo $t$ dipende da valori del segnale di ingresso in istanti $t' \neq t$.
Ad esempio, vediamo che l'integrale stesso è un operatore, che quindi rappresenta un sistema:
$$
y(t) = \int_{-\infty}^t x(\alpha) \, d\alpha
$$
dove i valori di $y(t)$ dipendono da $x(t)$ in istanti $\alpha < t$. Possiamo quindi scrivere anche:
$$
y(t) = \mathcal{T} \left[ x(\alpha) \ (\alpha < t); t \right]
$$

\subsubsection{Proprietà dei sistemi}
Vediamo in rassegna alcune possibili proprietà che possono avere sistemi.

\begin{enumerate}
	\item \begin{theorem}{Linearità dei sistemi}
		Un sistema si dice lineare se vale:
		$$
		\mathcal{T} \left[ a_1 x_1(\alpha) + a_2 x_2(\alpha); t \right] = a_1 y_1(t) + a_2 y_2(t)
		$$
	\end{theorem}
	questo, chiaramente, dato:
	$$
	\begin{cases}
		\mathcal{T} \left[ a_1 x_1(\alpha) \right] = a_1 y_1(t) \\
		\mathcal{T} \left[ a_2 x_2(\alpha) \right] = a_2 y_2(t) \\
	\end{cases}
	$$
	\item \begin{theorem}{Tempo invarianza dei sistemi}
		Un sistema è tempo invariante o \textit{stazionario} quando le sue proprietà non variano nel tempo, cioè:
		$$
		y(t - t_0) = \mathcal{T} \left[ x(\alpha - t_0); t \right]
		$$
	\end{theorem}
	Solitamente, un sistema è tempo variante quando la $t$ compare nell'equazione del sistema stesso, fuori da $x(\alpha)$;
	\item \begin{theorem}{Causalità dei sistemi}
		Un sistema si dice causale quando la sua uscita dipende solo da valori passati dell'ingresso, cioè: 
		$$
		y(t) = \mathcal{T} \left[ x(\alpha) \ (\alpha < t); t \right]
		$$
	\end{theorem}
	Un'esempio è l'integratore che abbiamo visto prima;
	\item \begin{theorem}{Assenza di memoria dei sistemi}
		Un sistema si dice senza memoria quando la sua uscita dipende solo dall'ingresso allo stesso istante, cioè: 
		$$
		y(t) = \mathcal{T} \left[ x(\alpha) \ (\alpha = t); t \right]
		$$
	\end{theorem}
\end{enumerate}

\subsubsection{Sistemi lineari stazionari}
Consideriamo quindi una particolare categoria di sistemi lineari, rappresentata dai \textit{Sistemi Lineari Stazionari} (\textbf{SLS}). Questi godono sia della proprietà di linearita che di stazionarietà. La loro uscita si ottiene come convoluzione:
$$
y(t) = x(t) \circledast h(t) = \int x(\beta) h(t - \beta) \, d\beta
$$
dove $h(t)$ è la risposta impulsiva del sistema, cioè alla delta di Dirac:
$$
h(t) = \mathcal{T} \left[ \delta(\alpha); t \right]
$$

Questo si giustifica vedendo che la convoluzione:
$$
x(\alpha) \circledast \delta(\alpha) = x(\alpha)
$$
è uguale ad $x$ valutata in $\alpha$ stessa. Svolgendo si ottiene quindi:
$$
y(t) = \mathcal{T} \left[ x(\alpha) ; t \right] = \mathcal{T} \left[ x(\alpha) \circledast \delta(\alpha) ; t \right] =
\mathcal{T} \left[ \int x(\beta) \delta(\alpha - \beta) \, d\beta ; t \right]
$$
Applicando la linearità dell'integrale sul sistema si ottiene:
$$
= \int \mathcal{T} \left[ x(\beta) \delta(\alpha - \beta) \right] \, d\beta = 
\int x(\beta) \mathcal{T} \left[ \delta(\alpha - \beta) \right] \, d\beta
$$
A questo punto si ha che, per la stazionarietà del sistema, vale rispetto alla risposta impulsiva:
$$
h(t) = \mathcal{T} \left[ \delta(\alpha); t \right] \implies
\mathcal{T} \left[ \delta(\alpha - \beta) \right] = h(t - \beta)
$$
per cui si ottiene il risultato finale:
$$
y(t) = \int x(\beta) h(t - \beta) \, d\beta = x(t) \circledast h(t)
$$

In qualche modo abbiamo quindi che la risposta impulsiva del sistema SLS rappresenta tutta la risposta del SLS, in quanto facendo la convoluzione con un segnale $x(t)$ di ingresso arbitrario si ottiene il corrispettivo segnale di uscita $y(t)$.

\par\smallskip

Vediamo l'esempio del sistema integratore:
$$
y(t) = \int_{-\infty}^t x(\alpha) \, d\alpha
$$
In questo caso la risposta impulsiva sarà data dall'integrale della delta di Dirac, ovvero dal gradino (per definizione della delta di Dirac stessa):
$$
h(t) = \int_{-\infty}^t \delta(\alpha) \, d\alpha = u(t)
$$
Questo si poteva dimostrare anche espandendo la convoluzione:
$$
y(t) = x(t) \circledast h(t) = \int_{-\infty}^{\infty} x(\alpha) h(t - \alpha) \, d\alpha
$$
che corrisponde all'integratore solo nel caso in cui $h(t) = u(t)$, in quanto:
$$
\int_{-\infty}^{\infty} x(\alpha) u(t - \alpha) \, d\alpha = \int_{-\infty}^t x(\alpha) \, d\alpha = y(t)
$$

\par\smallskip

\noindent
\textbf{\textsf{Risposta in frequenza}}

\noindent
Applicando il teorema della convoluzione ad un sistema SLS, si ottiene il risultato nel dominio frequenziale:
$$
y(t) = x(t) \circledast h(t) \implies Y(f) = H(f) X(f)
$$
questo direttamente dal teorema della convoluzione, e con $H(f)$ trasformata di $h(t)$. Chiamiamo questa trasformata della risposta impulsiva \textit{risposta in frequenza} del sistema SLS. Notiamo quindi che, in dominio frequenziale, la trasformata dell'uscita si ottiene attraverso una semplice operazione algebrica fra la trasformata dell'ingresso e la risposta in frequenza del sistema.

\par\smallskip

Vediamo ad esempio che, riguardo al sistema integratore, si ha che la risposta in frequenza è:
$$
u(t) = \frac{1}{2} + \frac{1}{2} \mathrm{sgn}(t) \rightarrow \frac{1}{2} \delta(t) + \frac{1}{j 2 \pi f}
$$
come calcolato alla fine della lezione 11. Questo ci permette ad esempio di dare una dimostrazione del teorema di \textit{integrazione} della trasformata di Fourier, che ci dice che la trasformata dell'integrale è:
$$
\mathcal{F} \left[ \int_{-\infty}^t x(\alpha) \, d\alpha \right] = X(f) H(f) = \frac{1}{2} X(0) + \frac{X(f)}{j 2 \pi f}
$$
Osserviamo come il termine $\frac{1}{2}\delta(t)$ nella trasformata del gradino, apparentemente arbitrario, ci fornisce semplicemente il termine condizione iniziale di un integrale trasformato. Abbiamo quindi che nel dominio frequenziale l'operazione abbastanza complicata dell'integrazione si traduce in una semplice divisione algebrica

\par\smallskip

Allo stesso modo possiamo dimostrare il teorema di derivazione, semplicemente invertendo quanto appena trovato:
$$
\frac{d x(t)}{dt} = X(f) \cdot j 2 \pi f
$$
dove chiaramente scompare il termine condizione iniziale (la derivata non ha memoria), e la divisione diventa una moltiplicazione.

\par\smallskip

\noindent
\textbf{\textsf{Sistemi in cascata}}

\noindent
Nel caso di più sistemi $\mathcal{T}_1$, $\mathcal{T}_2$ in cascata, cioè dove l'uscita $z(t)$ di un primo sistema viene data in ingresso ad un secondo sistema, si ottiene un un unico sistema la cui risposta impulsiva è data da:
$$
h(t) = h_1(t) \circledast h_2(t)
$$
cioè basta fare la convoluzione, da cui segue che la risposta in frequenza è:
$$
H(f) = H_1(f) H_2(f)
$$

\subsection{Filtri}
Un sistema lineare può rappresentare un \textbf{filtro}. In questo, un filtro non è altro che un sistema appositamente ricavato per eliminare alcune regioni di frequenze nel dominio frequenziale, ed amplificarne altre. Vediamone alcuni esempi:

\begin{itemize}
\item \textbf{\textsf{Filtro passa-basso ideale}} \\
Un filtro passa-basso ideale lascia passare inalterate le componenti a bassa frequenza (al di sotto di una certa frequenza di soglia detta di \textit{cutoff}), tagliando le altre. Lo abbiamo già introdotto in 7.2, e lo esprimiamo come la funzione $\mathrm{rect}$ nel dominio frequenziale:
$$
H(f) = \mathrm{rect} \left( \frac{f}{2B} \right)
$$
dove $B$ è la \textbf{banda} del filtro (per noi la nostra frequenza di cutoff).
Una conseguenza di questa scelta per la risposta in frequenza è che la risposta all'impulso (in dominio temporale) sarà la funzione $\mathrm{sinc}$:
$$
h(t) = 2 B \mathrm{sinc}(2B t)
$$
Questo è in accordo con quanto avevamo detto in 7.2, cioè che uno spettro di frequenze limitate da un filtro passa basso ideale si ricostruisce in dominio temporale come combinazione lineare di componenti $\mathrm{sinc}$.

\item \textbf{\textsf{Filtro passa-alto ideale}} \\
Dopo aver realizzato un filtro passa-basso ideale, è semplice realizzare un passa-alto. In dominio frequenziale, infatti, questo sarà semplicemente:
$$
H(f) = 1 - \mathrm{rect} \left( \frac{f}{2B} \right)
$$
Diventa quindi interessante vedere la risposta in dominio temporale di questo filtro, che non sarebbe stata affatto immediata senza l'intuizione che otteniamo nel dominio frequenziale. Questa sarà quindi:
$$
H(f) \rightarrow \delta(t) - 2 B \mathrm{sinc}(2B t)
$$

\item \textbf{\textsf{Filtro passa-banda ideale}} \\
Per la realizzazione del filtro passa-banda possiamo sfruttare il teorema di modulazione (introdotto sempre in 7.2) per spostare un filtro passa-basso (simmetrico rispetto all'origine degli assi delle frequenze) di una certa frequenza centrale $f_0$. In questo caso poniamo quindi:
$$
H(f) = \mathrm{rect} \left( \frac{f - f_0}{B} \right) + \mathrm{rect} \left( \frac{f + f_0}{B} \right)
$$
dove stavolta si è presa la banda $B$ come l'intera larghezza della banda che viene fatta passare attorno a $f_0$.
In dominio temporale questo si riporta alla modulazione per il termine $\cos (2 \pi f_0 t)$:
$$
h(t) = B \mathrm{sinc} (Bt) \cos(2 \pi f_0 t)
$$
\item \textbf{\textsf{Filtro elimina-banda ideale}} \\
Infine, mettiamo insieme quanto visto per il filtro passa-alto e per il fisso passa-banda per invertire il secondo ottenendo un filtro elimina-banda. In questo caso poniamo:
$$
H(f) = 1 - \mathrm{rect} \left( \frac{f - f_0}{B} \right) + \mathrm{rect} \left( \frac{f + f_0}{B} \right)
$$
dove ancara si è presa la banda $B$ come l'intera larghezza della banda che viene eliminata attorno a $f_0$.
In dominio temporale questo si riporta alla modulazione per il termine $\cos (2 \pi f_0 t)$, sottratta alla delta:
$$
h(t) = \delta(t) - B \mathrm{sinc} (Bt) \cos(2 \pi f_0 t)
$$
\end{itemize}

\subsubsection{Distorsione nei filtri}
Un filtro viene detto \textbf{non distorcente} se nell'intorno della sua banda lo si può scrivere come:
$$
H(f) = K \cdot e^{-j 2 \pi f t_0}
$$
cioè nella banda che viene lasciata passare:
\begin{itemize}
	\item La risposta in ampiezza è \textit{piatta};
	\item La risposta in fase è \textit{lineare} (è impossibile realizzare un filtro che la mantenga sempre piatta, a meno del caso passa-tutto).
\end{itemize}

Al di fuori della banda che viene fatta passare, chiaramente, il filtro può assumere qualsiasi forma (non potrà esseree piatto nelle regioni che inizia a tagliare).

\end{document}
